\documentclass[12pt]{article}
\usepackage[T1]{fontenc}
\usepackage[utf8]{inputenc}
\usepackage{lmodern}
\usepackage{textcomp}
\usepackage{enumitem}
\usepackage{geometry}
\geometry{margin=1in}
\begin{document}
\begin{center}
Answer Key - Religious Studies - K-12 Level Assessment 9 \
\small (Answers are ordered exactly as in the question paper.)
\end{center}

\section*{Multiple Choice}
\begin{enumerate}[label=\arabic*.]
\item \textbf{Answer:} B
\\ \textit{Options:} A) Second century CE; B) Fifth century CE; C) Fifth century BCE; D) Eighteenth century CE
\\ \textit{Note:} The first Jaina temples appeared in the fifth century CE as this period marked significant developments in Jaina architecture and the establishment of dedicated places of worship reflective of the religion's growing prominence in Indian society.
\item \textbf{Answer:} C
\\ \textit{Options:} A) Khadijah; B) Gabriel; C) Waraqah; D) Abu Bakr
\\ \textit{Note:} Waraqah, a cousin of Khadijah, recognized Muhammad's prophetic calling after he received his first revelation, affirming him as a prophet in the early days of Islam.
\item \textbf{Answer:} A
\\ \textit{Options:} A) 48; B) 60; C) 12; D) 36
\\ \textit{Note:} Jaina devotees practice meditation or reflection for 48 minutes each day as part of their spiritual discipline to cultivate mindfulness and inner peace.
\item \textbf{Answer:} B
\\ \textit{Options:} A) The rainbow; B) Circumcision; C) A son; D) Bar mitzvah
\\ \textit{Note:} Circumcision is the sign of the covenant for Jewish males as it is a key ritual established in the Torah that symbolizes their special relationship with God and the commitment to follow His commandments.
\item \textbf{Answer:} C
\\ \textit{Options:} A) Parvati; B) Sita; C) Sri (Lakshmi); D) Sarasvati
\\ \textit{Note:} Sri, also known as Lakshmi, is the goddess of good fortune and prosperity and is traditionally depicted as the consort of Vishnu in many of his incarnations.
\end{enumerate}

\section*{True / False}
\begin{enumerate}[label=\arabic*.]
\setcounter{enumi}{5}
\item \textbf{Answer:} False
\\ \textit{Options:} A) True; B) False
\\ \textit{Note:} In Sikhism, the communal meal offered at the place of worship is referred to as Langar, not Sangat, which actually denotes the congregation or community of Sikhs gathering for worship.
\item \textbf{Answer:} False
\\ \textit{Options:} A) True; B) False
\\ \textit{Note:} Zen Buddhism encourages lay participation in its community activities, emphasizing the importance of integrating meditation and mindfulness into everyday life rather than isolating it to monastic practice.
\item \textbf{Answer:} True
\\ \textit{Options:} A) True; B) False
\\ \textit{Note:} The Daodejing is indeed referred to in English as the Classic of the Way and Power, reflecting its central themes of the Dao (Way) and De (Power or Virtue) in Daoist philosophy.
\item \textbf{Answer:} False
\\ \textit{Options:} A) True; B) False
\\ \textit{Note:} The statement is false because the Egyptian goddess Hathor was primarily depicted as a cow or as a woman with cow horns, symbolizing fertility and motherhood, rather than as an eagle.
\item \textbf{Answer:} True
\\ \textit{Options:} A) True; B) False
\\ \textit{Note:} Ojizo-sama is revered in Japanese Buddhism as a bodhisattva who protects and guides the souls of deceased children, helping them find peace and salvation in the afterlife.
\end{enumerate}

\section*{Long Form Response}
\begin{enumerate}[label=\arabic*.]
\setcounter{enumi}{10}
\item \textbf{Answer:} In Sikhism, remembering the Divine Name, or "Naam," is of great significance as it serves as a means to connect with God and attain spiritual liberation. This practice involves reciting and meditating on God's name, which helps individuals cultivate a sense of inner peace and awareness of the divine presence in their lives. By focusing on the Divine Name, Sikhs believe they can overcome worldly distractions and ego, leading to a deeper understanding of their spiritual purpose. Ultimately, this remembrance fosters a relationship with God that is essential for achieving liberation from the cycle of birth and rebirth, allowing individuals to experience unity with the divine.
\\ \textit{Note:} This answer is correct because it emphasizes that remembering the Divine Name in Sikhism is essential for establishing a connection with God, fostering inner peace, and overcoming ego, all of which contribute to spiritual liberation.
\item \textbf{Answer:} Mahavira holds significant importance in Jaina traditions as the last of the ascetic prophets. His teachings emphasized the principles of non-violence (ahimsa), truth, and self-discipline, which became fundamental tenets of Jainism. By promoting a strict ascetic lifestyle, Mahavira inspired his followers to pursue spiritual purity and liberation from the cycle of rebirth. His emphasis on personal responsibility and ethical conduct profoundly shaped the faith, leading to a strong community dedicated to these principles. Overall, Mahavira's contributions not only defined Jaina beliefs but also established a lasting framework for spiritual practice within the religion.
\\ \textit{Note:} correct because it accurately highlights Mahavira's role as the last ascetic prophet in Jainism and emphasizes his teachings on non-violence, truth, and self-discipline, which are foundational to the faith and have significantly influenced the ethical and spiritual practices of his followers.
\end{enumerate}

\vfill
\noindent\textit{End of Answer Key}
\end{document}